\documentclass{situationist-preprint}

\title{Metaphor and Literal Falsity}
\msid{LNG-09-b}
\msfamily{LNG-09}
\msdiscipline{Philosophy of Language}
\mstype{Consequence}
\mskeywords{metaphor; literal meaning; truth; Davidson; implicature; assessment}

\begin{document}
\maketitle

\begin{abstract}
If metaphors mean exactly what they literally mean, as LNG-09-a argues, then the truth value of a metaphorical utterance is the truth value of its literal content. We draw three consequences. Most metaphors are false, and their falsity is usually patent; some are true, so literal falsity is neither necessary nor sufficient for metaphor. A speaker who utters a false metaphor neither lies nor errs, since she does not put the literal content forward as believed. And when we call a metaphor \enquote{true}, we assess its use, what it conveys or makes us notice, not its content. We show that this assessment behaves unlike truth in inference and embedding.
\end{abstract}

\section{The truth value of a metaphor}

Theorem~3.3 of LNG-09-a holds that metaphors mean exactly what they literally mean, and that their distinctive effect lies in use. Since a sentence's truth value on a use is fixed by what it means on that use, we have:

\begin{proposition}\label{p:tv}
A metaphorical utterance is true if and only if its literal content is true.
\end{proposition}

\begin{corollary}\label{c:false}
\enquote{Juliet is the sun} and \enquote{Man is a wolf to man}, uttered metaphorically, are false.
\end{corollary}

Most metaphors are patently false, and Davidson observed that this is often how we recognize them \citep{Davidson1978}. Grice classed metaphor with irony as a flouting of the first maxim of Quality: the speaker says what both know to be false, and the hearer, assuming cooperation, looks for what else she might be conveying \citep{Grice1975}. On that account the literal falsity is the trigger for the search.

But literal falsity is not necessary for metaphor. \enquote{No man is an island} is literally true, and it is used metaphorically; Davidson made the same point with this example \citep{Davidson1978}. Nor is it sufficient: an ordinary mistake is literally false and not metaphorical. The controversion theory, on which a metaphor is a sentence whose literal attribution is logically or obviously false \citep{Beardsley1962}, therefore fails as a definition. What makes an utterance metaphorical is that the speaker uses its literal content to get the hearer to see one thing as another; falsity is the commonest means of signaling this, not its essence.

\section{Neither lie nor error}

\begin{proposition}\label{p:nolie}
A speaker who utters a literally false metaphor neither lies nor makes a mistake.
\end{proposition}

\begin{proof}[Argument]
A lie requires that the speaker assert what she believes false, and a mistake requires that she believe what she asserts. In uttering a patently false metaphor the speaker does not put the literal content forward as something she believes: it is mutually evident that she does not, as it is in irony. Since she does not assert the content in the sense in which lies and mistakes involve assertion, she neither lies nor errs. What she does commit herself to, if anything, is what she implicates, which is not the literal content.
\end{proof}

This explains why the falsity of most metaphors, which is a consequence of Davidson's thesis, does not make metaphorical speech defective. The objection that Davidson's thesis convicts poets of uttering falsehoods confuses falsity of content with defect of use.

\section{Calling a metaphor true}

We often say that a metaphor is true, apt or deep. Bergmann argued that a metaphorical utterance is an assertion with a truth value of its own, fixed by what the speaker means \citep{Bergmann1982}. If so, \enquote{true} in \enquote{that metaphor is true} would ascribe truth to a non-literal content, contrary to Proposition~\ref{p:tv}. We argue that it does not.

\begin{proposition}\label{p:assess}
In \enquote{what she said is true}, said of a metaphorical utterance whose literal content is false, \enquote{true} assesses the use of the utterance, not its content.
\end{proposition}

\begin{proof}[Argument]
If \enquote{true} here ascribed truth to a content, it would license the inferences truth licenses. From \enquote{Juliet is the sun} and \enquote{what Romeo said is true}, one could infer that Juliet is a star, which no competent speaker does; from the same premises, one could disquote, and no one does. The assessment does license other inferences: that Romeo's metaphor brings out something real about Juliet, or about his love for her. These are claims about what the utterance conveys or makes salient, and so about its use. By Principle~P2 of LNG-09-a, what a speaker brings about by uttering a sentence, including implicatures and changes of attention, belongs to use.
\end{proof}

The assessment is of the same kind as calling a remark apt, a joke good or an analogy illuminating. It can be correct or incorrect; but its correctness does not consist in the truth of a proposition the sentence expresses. Where the metaphor carries a determinate implicature, as \enquote{Man is a wolf to man} implicates that people are often ruthless toward one another, calling it true may also be a way of endorsing the implicature, which is a proposition; but the implicature is calculable from what was said and cancellable (LNG-11), and so is not what the sentence means.

\begin{objection}
Disquotation is not a test, since \enquote{true} applied to any utterance can be used loosely.
\end{objection}

\begin{reply}
The test is not that loose uses of \enquote{true} are impossible but that this use is systematically non-disquotational. For literal utterances, assessments of truth and the corresponding first-order claims stand or fall together; for metaphors they come apart in every case. That pattern is what one expects if \enquote{true} assesses a different object.
\end{reply}

\begin{objection}
In \enquote{If Juliet is the sun, Romeo's night is over}, the metaphor embeds, and the conditional seems to have a truth value that depends on a metaphorical content.
\end{objection}

\begin{reply}
The embedded metaphor is interpreted as a whole: the speaker uses an extended image, and what she conveys is again a matter of use. The literal conditional has the truth value its literal parts give it; the impression that it says something true about Romeo derives from the same source as the impression that the simple metaphor does.
\end{reply}

\section{Conclusion}

From the thesis that metaphors mean what they literally mean, it follows that most metaphors are false and some true, that their falsity is no defect in the speaker, and that calling a metaphor true assesses its use, not its content.

\begin{thebibliography}{9}\small
\bibitem[Beardsley(1962)]{Beardsley1962} Beardsley, M.~C. (1962). The metaphorical twist. \emph{Philosophy and Phenomenological Research}, 22, 293--307.
\bibitem[Bergmann(1982)]{Bergmann1982} Bergmann, M. (1982). Metaphorical assertions. \emph{Philosophical Review}, 91, 229--245.
\bibitem[Davidson(1978)]{Davidson1978} Davidson, D. (1978). What metaphors mean. \emph{Critical Inquiry}, 5, 31--47.
\bibitem[Grice(1975)]{Grice1975} Grice, H.~P. (1975). Logic and conversation. In P.~Cole and J.~L. Morgan (eds.), \emph{Syntax and Semantics 3: Speech Acts}. Academic Press.
\end{thebibliography}

\end{document}
